\subsection{Actual finite-fibre multiplicity conservation under base change}
\label{subsec:finite-fibre-multiplicity-base-change}

All pointwise results below need only module finiteness of the original
commutative-ring algebra. A base-changed point is given; existence of points
above every base point is not asserted. Every fibre uses its own genuine
local length times residue degree. Those weights are conserved by sums over
restriction fibres, rather than by a bijection of the individual points.

\begin{lemma}[Actual local-factor idempotents and their entire norm forms]
\label{algebra:finite-factor-idempotent-norm}
\lean{EnvelopeProperty.normForm_algebraMap,
EnvelopeProperty.finiteAlgebraIdempotent,
EnvelopeProperty.finiteAlgebraIdempotent_algebraMap,
EnvelopeProperty.finiteAlgebraIdempotent_isIdempotent,
EnvelopeProperty.finiteAlgebraIdempotent_mem_iff,
EnvelopeProperty.localization_idempotent_algebraMap,
EnvelopeProperty.normForm_idempotent}
\leanok
\uses{algebra:finite-artinian-local-factors,spectrum:whole-norm-form-base-change}
Let $E$ be a finite commutative algebra over any field $H$. The actual
localization product constructs an idempotent $e_m$ equal to one on $E_m$
and zero on every other local factor. It belongs to a maximal ideal $p$
exactly when $p\ne m$. For any idempotent $e\in E$,
\[
 N_{E/H}(Te+U(1-e))=
 \prod_{p\in\operatorname{Max}(E)}
 \begin{cases}T^{\dim_HE_p},&e\notin p,\\
 U^{\dim_HE_p},&e\in p.
 \end{cases}
\]
These are equalities of the whole determinant polynomials.
\end{lemma}
\begin{proof}
Invert the genuine Artinian localization product on the coordinate vector
with value one at $m$ and zero elsewhere. Its coordinate formula proves
idempotence. Localization membership in its maximal ideal detects original
membership, yielding the stated criterion. In an arbitrary local ring an
idempotent is zero or one: either it or its complement is a unit, and
unit cancellation in $e(1-e)=0$ gives the result.

The real product isomorphism factors the whole norm form. In a local factor
$e$ is the scalar zero or one, so its linear form is $U$ or $T$.
For a scalar family, the balancing identity in the actual polynomial tensor
algebra makes the linear form a scalar coefficient-ring element; its scalar
norm is that polynomial raised to the actual free rank. Applying this to all
local factors gives the formula. No field-point interpolation or reduction
of nilpotent coefficients occurs.
\end{proof}

\begin{theorem}[Genuine local dimension conservation for every field extension]
\label{algebra:local-dimension-base-change}
\lean{EnvelopeProperty.degreeOf_map_normForm_idempotent,
EnvelopeProperty.degreeOf_normForm_idempotent,
EnvelopeProperty.finiteAlgebraMaximalComap,
EnvelopeProperty.finiteAlgebraIdempotent_baseChange_not_mem_iff,
EnvelopeProperty.sum_localFinrank_baseChange}
\leanok
\uses{algebra:finite-factor-idempotent-norm,spectrum:whole-norm-form-base-change}
For any field extension $K/H$, set $F=K\otimes_HE$. Contraction along the
actual right inclusion maps each maximal ideal $p$ of $F$ to a maximal
ideal of $E$. For every maximal ideal $m$ of $E$,
\[
 \sum_{p\cap E=m}\dim_KF_p=\dim_HE_m.
\]
No valuation, completeness, separability or reducedness is required.
\end{theorem}
\begin{proof}
Contraction is prime, and $E$ is Artinian, so it is genuinely maximal.
Apply whole norm-polynomial coefficient base change to
$Te_m+U(1-e_m)$. Compute degree in $T$ using the preceding factorization,
where degrees of nonzero polynomial products and powers add over a field.
Before extension, only the factor $m$ contributes, with degree $\dim_HE_m$.
After extension, a maximal factor contributes its whole dimension precisely
when $1\otimes e_m$ is outside its maximal ideal. The original idempotent
membership criterion identifies this condition with contraction equal to
$m$. Equality of the actual polynomials gives the dimension identity.
The complete local dimensions, including all nilpotent and inseparable
thickness, remain in the degree calculation.
\end{proof}

\begin{theorem}[Arbitrary function products on maximal factors]
\label{algebra:finite-factor-weighted-products}
\lean{EnvelopeProperty.finiteAlgebraWeightedProduct_baseChange}
\leanok
\uses{algebra:local-dimension-base-change}
For every commutative monoid $M$ and every function
$a:\operatorname{Max}(E)\to M$,
\[
 \prod_{p\in\operatorname{Max}(F)}a(p\cap E)^{\dim_KF_p}
 =\prod_{m\in\operatorname{Max}(E)}a(m)^{\dim_HE_m}.
\]
\end{theorem}
\begin{proof}
Regroup the finite product by the actual contraction map. Within each
contraction fibre the base element $a(m)$ is fixed, so its powers multiply
by adding their exponents. The preceding genuine dimension conservation
replaces that sum by $\dim_HE_m$. Empty spectra give the empty product.
Neither injectivity nor a bijection of contraction fibres is assumed.
\end{proof}

\begin{lemma}[Actual maximal localization comparison under an algebra isomorphism]
\label{algebra:maximal-localization-isomorphism}
\lean{EnvelopeProperty.maximalSpectrumAlgEquiv,
EnvelopeProperty.maximalLocalizationAlgEquiv,
EnvelopeProperty.finrank_maximalLocalizationAlgEquiv}
\leanok
An actual coefficient-algebra isomorphism $E\simeq F$ induces a bijection
of maximal ideals and actual coefficient-algebra isomorphisms
$E_{p\cap E}\simeq F_p$. Over a field these preserve local dimensions.
\end{lemma}
\begin{proof}
Pull ideals back through the ring isomorphism and its inverse. Both
compositions are identity on every ideal, and maximality is preserved.
The isomorphism maps the actual complements of these primes onto each
other. The genuine localization equivalence for this mapped submonoid is
therefore a coefficient-algebra isomorphism. Its underlying linear
equivalence preserves dimension. This construction works over arbitrary
commutative coefficient rings, before the field dimension specialization.
\end{proof}

\begin{definition}[Genuine completed fibre tensor base change]
\label{spectrum:completed-fibre-base-change-tensor}
\lean{EnvelopeProperty.Berkovich.Spectrum.fibreBaseChangeEquiv,
EnvelopeProperty.Berkovich.Spectrum.fibreBaseChangeEquiv_one_tmul,
EnvelopeProperty.Berkovich.Spectrum.baseChangeFibreRestrict,
EnvelopeProperty.Berkovich.Spectrum.fibreBaseChangeEquiv_symm_tmul}
\leanok
\uses{spectrum:fibre-tensor-evaluation,spectrum:completed-residue-field-pullback}
Let $A\to B$ and $A\to C$ be arbitrary commutative-ring algebra maps,
let $x'\in\mathcal S(C)$, and let $x$ be its restriction to $A$. Actual
completed-field maps give the scalar action $H(x)\to H(x')$ and the genuine
$H(x')$-algebra isomorphism
\[
 H(x')\otimes_C(C\otimes_AB)
 \simeq H(x')\otimes_{H(x)}(H(x)\otimes_AB).
\]
It preserves every original algebra element. Restriction along $B\to
C\otimes_AB$ maps the actual fibre over $x'$ into the actual fibre over $x$.
\end{definition}
\begin{proof}
The real isometric completed-field map and its all-coefficient commuting
square define the required scalar actions and towers. Two actual tensor
cancellation isomorphisms identify both displayed algebras with
$H(x')\otimes_AB$. Their composite proves the asserted isomorphism and
its pure-tensor formulas. Balancing the original coefficient $a\in A$
inside $C\otimes_AB$ shows that every restricted point has base value
$x(a)$, proving the genuine fibre restriction.
\end{proof}

\begin{theorem}[Whole fibre tensor maps commute with restriction]
\label{spectrum:completed-fibre-base-change-naturality}
\lean{EnvelopeProperty.Berkovich.Spectrum.completedMap_baseChange_comm,
EnvelopeProperty.Berkovich.Spectrum.fibreHom_baseChange_old}
\leanok
\uses{spectrum:completed-fibre-base-change-tensor,spectrum:completed-residue-field-pullback}
For a point $y'$ over $x'$ restricting to $y$ over $x$, the new completed
tensor map restricted through the displayed tensor comparison and right
inclusion agrees on every tensor with the old completed tensor map followed
by the genuine isometric map $H(y)\to H(y')$.
\end{theorem}
\begin{proof}
Actual completed-field composition and the balanced coefficient square
prove equality of the two maps $H(x)\to H(y')$. The tensor comparison
formula then proves the asserted equality on each pure tensor: both maps
are the same completed scalar value times the same original algebra
value in $H(y')$. Tensor induction on zero, pure tensors and sums proves
it on the whole algebra. This uses genuine completion-map functoriality,
not equality of numerical values on a finite generator family.
\end{proof}

\begin{theorem}[Actual restricted maximal kernels and local multiplicities]
\label{spectrum:completed-fibre-base-change-support}
\lean{EnvelopeProperty.Berkovich.Spectrum.baseChangeFibreMaximal,
EnvelopeProperty.Berkovich.Spectrum.baseChangeFibreMaximal_comap,
EnvelopeProperty.Berkovich.Spectrum.fibreNormMult_eq_localFinrank_of_comap,
EnvelopeProperty.Berkovich.Spectrum.fibreNormMult_baseChange_eq_dim}
\leanok
\uses{spectrum:completed-fibre-base-change-naturality,
algebra:maximal-localization-isomorphism,spectrum:finite-completed-fibres,
spectrum:nonetale-fibre-multiplicities}
Assume only that $B$ is finite as an $A$-module. The actual maximal
completed tensor kernel of $y'$ transports through the tensor isomorphism
to a maximal ideal whose contraction is exactly the actual kernel of $y$.
Its local $H(x')$-dimension is precisely the genuine norm multiplicity of
$y'$. The old local $H(x)$-dimension is the genuine norm multiplicity of $y$.
\end{theorem}
\begin{proof}
Finiteness makes both completed tensor supports maximal by the genuine
integral fibre classification. Transport the new ideal through the actual
algebra isomorphism. The preceding whole-tensor map equality identifies its
contraction with the kernel of the old map followed by $H(y)\to H(y')$.
That field map is genuinely injective, so the kernel is exactly the old
one. The finite local-factor dimension identity is the genuine local
length times residue degree. Actual localization equivalence under the
tensor isomorphism preserves the new local dimension, giving the claimed
multiplicities at named base points.
\end{proof}

\begin{theorem}[All function products and sums under finite base change]
\label{spectrum:all-function-finite-base-change}
\lean{EnvelopeProperty.Berkovich.Spectrum.fibreWeightedProduct_baseChange,
EnvelopeProperty.Berkovich.Spectrum.fibreWeightedSum_baseChange,
EnvelopeProperty.Berkovich.Spectrum.fibreWeightedERealSum_baseChange}
\leanok
\uses{algebra:finite-factor-weighted-products,
spectrum:completed-fibre-base-change-support}
For $B/A$ finite, the given $x'$ and any function $a:\mathcal S(B)\to M$
to a commutative monoid,
\[
 \prod_{y'\text{ over }x'}a(y'|_B)^{m_{y'}}
 =\prod_{y\text{ over }x}a(y)^{m_y}.
\]
For every function to an additive commutative monoid the analogous weighted
sum identity holds. In particular, for every
$u:\mathcal S(B)\to\overline{\mathbb R}$,
\[
 \sum_{y'\text{ over }x'}m_{y'}u(y'|_B)
 =\sum_{y\text{ over }x}m_yu(y).
\]
No freeness, projectivity, flatness, regularizability or upper-finiteness
of $u$ is needed.
\end{theorem}
\begin{proof}
The genuine maximal-kernel fibre equivalences reindex the new and old
finite point products by the maximal ideals of their actual completed
tensor algebras. The whole support comparison identifies restriction with
contraction after the tensor isomorphism, and the local dimension comparison
identifies every exponent. Apply the arbitrary-function algebraic product
identity and reindex back through the genuine old fibre equivalence.

Apply this product equality to the multiplicative type tag of an arbitrary
additive commutative monoid, then unwrap the tag. This gives the sum identity
with natural scalar multiples. Mathlib's EReal is an additive commutative
monoid, and its unconditional identity $n\mathbin{\bullet}t=(n:t)t$ turns
that sum into the displayed EReal multiplication formula. It uses Mathlib's
actual extended-real addition and zero-times-infinity conventions, including
mixed infinities. No additional finite-value restriction is introduced.
\end{proof}

\begin{corollary}[Multiplicity mass over each restricted point]
\label{spectrum:finite-base-change-multiplicity-mass}
\lean{EnvelopeProperty.Berkovich.Spectrum.sum_fibreNormMult_baseChange}
\leanok
\uses{spectrum:all-function-finite-base-change}
For every old point $y$ over $x$,
\[
 \sum_{\substack{y'\text{ over }x'\\y'|_B=y}}m_{y'}=m_y.
\]
\end{corollary}
\begin{proof}
Apply the additive function identity to the natural-valued indicator of
$y$. Its old weighted sum has a single nonzero term because $y$ belongs to
the old fibre. The new weighted sum selects exactly the points restricting
to $y$, giving the conservation law with all genuine local multiplicities.
\end{proof}
